% !TEX root = ../main.tex
\chapter{Sistemi di tipo}
\label{chap:type_systems}

L'introduzione dei tipi nel $\lambda$-calcolo risponde alla necessità di limitare l'espressività del formalismo puro in cambio di garanzie logiche e computazionali, 
prima fra tutte la terminazione del calcolo. In questo capitolo analizzeremo come le annotazioni sui termini permettano di prevenire errori semantici e di instaurare 
un controllo rigoroso sulla formazione delle espressioni. Il percorso inizierà con l'esame del $\lambda$-calcolo semplicemente tipato ($\lambda^{\to}$), di cui verranno 
definiti i processi di verifica (type checking) e di ricostruzione automatica (type inference). Successivamente, utilizzeremo il Cubo di Barendregt come cornice teorica 
per classificare le estensioni del calcolo lungo tre assi di astrazione: polimorfismo, operatori di tipo e tipi dipendenti. Il capitolo culminerà nell'analisi del Calcolo 
delle Costruzioni (CoC), dove l'unificazione delle categorie sintattiche permette di trattare termini, tipi e prove logiche all'interno di un unico, potente linguaggio formale.

\section{\texorpdfstring{$\lambda$}{Lambda}-calcolo semplicemente tipato (\texorpdfstring{$\lambda^{\to}$}{Lambda-to})}
Il $\lambda$-calcolo semplicemente tipato ($\lambda^{\to}$), come suggerisce il nome, si distingue dal sistema puro per l'introduzione dei tipi.
In questo sistema di riscrittura i tipi sono della annotazioni sulla varibile vincolata della astrazione

\subsection{Grammatica del \texorpdfstring{$\lambda$}{Lambda}-calcolo semplicemente tipato}
La grammatica BNF del $\lambda$-calcolo semplicemente tipato ($\lambda^{\to}$), estende il calcolo puro tramite
la categoria sintattica dedicata ai tipi, legando le categorie sintattice tramite il costrutto dell'astrazione 
\eqref{eq:simple_typed_lambda_grammar}.
\begin{equation}
    \begin{array}{lrll}
        \text{Types:}   & \tau, \sigma & ::= \iota \quad                & \text{(Tipo base / atomico)}      \\
                        &              & \mid \tau_1 \to \tau_2 \quad   & \text{(Tipo funzione)}            \\
        \text{Terms:}   & e            & ::= x \quad                    & \text{(Variabile)}                \\
                        &              & \mid \lambda x : \tau . e \quad& \text{(Astrazione)}               \\
                        &              & \mid e_1 \, e_2 \quad          & \text{(Applicazione)}
    \end{array}
    \tag{STLC}
    \label{eq:simple_typed_lambda_grammar}
\end{equation}
Si osservi come la sintassi dell'astrazione includa il tipo, cio garantiche che ogni variabile vincolata sia univovamete tipata 
all'interno del proprio scope. Inoltre è opportuno ribadire che la grammatica non fornisce semantica.

\subsection{Type checking}
Il giudizio di tipizzazione o \emph{type checking} ($\Gamma \vdash e : \tau$) è il processo d'\emph{inferenza} mediante il quale si determina se un termine appertiene a un determinato tipo sotto un insieme di assunzioni.
Tale processo opera in modo induttivo sulla grammatica dei termini, visitando l'intera struttura verificandone le componenti. Le applicazioni ricorsive delle regole per verificare il tipo finale,
garantendo la correttezza formale dell'intera computazione.
\begin{equation}
    \text{type}(\Gamma, e) = 
    \begin{cases} 
        \tau & \text{if } e = x \text{ e } (x : \tau) \in \Gamma \\
        \sigma \to \text{type}(\Gamma, x:\sigma, e') & \text{if } e = \lambda x : \sigma . e' \\
        \tau & \text{if } 
            \begin{cases} 
            e = e_1 e_2 \\
            \text{type}(\Gamma, e_1) = \sigma \to \tau \\
            \text{type}(\Gamma, e_2) = \sigma 
            \end{cases} \\
        \perp & \text{otherwise}
    \end{cases}
    \label{eq:simple_typed_lambda_type_cheking_algorithm}
\end{equation}
Il contesto di tipizzazione, denotato con $\Gamma$, è un insieme finito di assunzioni di tipo nella forma $x : \sigma$.
L'estensione del contesto, indicata con $\Gamma, x : \sigma$, rappresenta l'aggiunta di una nuova assunzione, con la clausola 
che se $x$ era già presente in $\Gamma$, la nuova assegnazione sovrascrive la precedente (\emph{meccanismo di shadowing}). Il contesto
può essere modellizzato come una funzione parziale. Il simbolo $\perp$ (leggasi bottom) invece rapresenta il fallimento dell'algoritmo
di verifica, questo in seguito rapresentera una specifica condizione sul termine, limitiamoci a dire che è mal tipato (\emph{ill-typed}).

\subsection{Type inference}
L'inferenza di tipo o \emph{type inference}  è il prcesso mediante il quale si ricostruisce automaticamente il tipo più generale di un termine,
utilizando al struttura sintattica e le relazioni tra componenti.

\begin{equation}
    \text{gen}(\Gamma, e) = 
    \begin{cases} 
        (\tau, \emptyset) & \text{if } e = x \land (x : \tau) \in \Gamma \\
        (\alpha \to \tau, \mathcal{C}) & \text{if } e = \lambda x. e' \land (\tau, \mathcal{C}) = \text{gen}(\Gamma, x : \alpha, e') \\
        (\beta, \mathcal{C}_{all}) & \text{if } e = e_1 e_2 \text{ where:} \\ & 
        \begin{cases} 
            (\tau_1, \mathcal{C}_1) = \text{gen}(\Gamma, e_1) \\
            (\tau_2, \mathcal{C}_2) = \text{gen}(\Gamma, e_2) \\
            \mathcal{C}_{all} = \mathcal{C}_1 \cup \mathcal{C}_2 \cup \{\tau_1 = \tau_2 \to \beta\}
        \end{cases}
    \end{cases}
    \label{eq:constraint_generation}
\end{equation}
Il risultato del processo d'inferenza di tipo è un insieme di equazioni $\mathcal{C}$ che rapresentano i vincoli tra tipi delle componenti del termine.
Se il sistema di equazioni dei vicoli di tipi ha soluzione allora è possibile dimostrare che il risoltato è il tipo più generale che rispetta tali vincoli 
(\emph{MGU}). L'insieme di equasioni $\mathcal{C}$ viene anche chiamato \emph{problema di unificazione del primo ordine}.
\begin{equation}
    \mathcal{U}(\mathcal{C}) = 
    \begin{cases} 
        id & \text{if } \mathcal{C} = \emptyset \\
        \mathcal{U}(\mathcal{C}') & \text{if } \mathcal{C} = \{\tau = \tau\} \cup \mathcal{C}' \\
        \mathcal{U}(\mathcal{C}'[\alpha \mapsto \tau]) \circ [\alpha \mapsto \tau] & \text{if } \mathcal{C} = \{\alpha = \tau\} \cup \mathcal{C}' \text{ e } \alpha \notin FV(\tau) \\
        \mathcal{U}(\mathcal{C}'[\alpha \mapsto \tau]) \circ [\alpha \mapsto \tau] & \text{if } \mathcal{C} = \{\tau = \alpha\} \cup \mathcal{C}' \text{ e } \alpha \notin FV(\tau) \\
        \mathcal{U}(\{\tau_1 = \sigma_1, \tau_2 = \sigma_2\} \cup \mathcal{C}') & \text{if }    \mathcal{C} = \{\tau_1 \to \tau_2 = \sigma_1 \to \sigma_2\} \cup \mathcal{C}' \\
        \perp & \text{otherwise (failed)}
    \end{cases}
    \label{eq:unification_algorithm}
\end{equation}
La risoluzione del sistema di vincoli poggia su principi logici intuitivi ma rigorosi: 
l'impossibilità di far coincidere tipi atomici distinti (assenza di collisioni tra costruttori) e 
il divieto di formulare definizioni ricorsive, per cui un tipo non può essere definito in funzione di se stesso (occurs check).
\begin{table}[h]
    \centering \small
    \begin{tabular}{lp{5cm}l}
        \hline
        \textbf{Errore} & \textbf{Descrizione}                   & \textbf{Esempio di Vincolo}  \\ \hline
        Clash Error     & Unificazione di costruttori diversi    & $Bool = \sigma \to \tau$     \\
        Occurs Check    & Definizione di tipi ricorsivi infiniti & $\alpha = \alpha \to \beta$  \\
        Type Mismatch   & Incongruenza tra tipi base differenti  & $Nat = Bool$                 \\ \hline
    \end{tabular}
    \caption{Fallimenti dell'algoritmo di unificazione.}
    \label{tab:unification_failures}
\end{table}

\subsection{Normalizzazione forte}
Il sistema $\lambda^{\to}$ gode della proprita di \emph{normalizzazione forte}, ovvero ogni termine ben tipato riduce in un numero finito di passi
in forma normale indipendetemente dalla strategia di riduzione utilizzata. Per una trattazione esaustiva della dimostrazione si rimanda a \cite{Girard1989} o
\cite{Hindley1997}.

\section{Cubo di Barendregt}
Una volta definiti i fondamenti del $\lambda^{\to}$-calcolo, risulta necessario introdurre un modello per orientarsi e classificare
i sistemi di tipo estensionali del $\lambda^{\to}$-calcolo. Tale struttura trova sua espressione del \emph{Cubo di Barendregt} \cite{barendregt1992}.
I sistemi analizzati nel cubo non sono mutualmente esclusivi, ma cumulativi. Si parte del vertice del $\lambda^{\to}$-calcolo e ogni
asse aggiunge potere espressivo l sistema dei tipi, riportate di seguito le direzioni.

\begin{itemize}
    \item \textbf{Asse $x$ (Tipi Dipendenti, \texorpdfstring{$\lambda P$}{LP}):}  
        Tipi che dipendono da termini. Il tipo di un oggetto può variare in base al valore di un altro (es. un vettore di lunghezza $n$).
    \item \textbf{Asse $y$ (Polimorfismo, \texorpdfstring{$\lambda 2$}{L2}):} 
        Termini che dipendono da tipi. Permette di scrivere funzioni universali che accettano un tipo come argomento (es. la funzione identità generica).
    \item \textbf{Asse $z$ (Operatori di Tipo, \texorpdfstring{$\lambda \underline{\omega}$}{L-omega}):}  Tipi che dipendono da tipi.
        Permette di definire funzioni a livello dei tipi, ovvero costruttori che generano nuovi tipi a partire da tipi esistenti. 
\end{itemize}
Il sistema di tipi opposto al vertice del $\lambda^{\to}$-calcolo, e dunque il più espressivo, prende il nome di \emph{Calcolo delle Costruzioni (CoC)}. 
In tale sistema, il confine tra l'esecuzione del codice (valutazione dei termini) e il type checking si perdono. Nei sistemi con tipi dipendeti, l'ugualianza 
tra due tipi puo dipendere dalla riduzione di un termine in forma normale. 

\subsection{Sistema \texorpdfstring{$\lambda 2$}{L2}}

\paragraph{Grammatica.}
La grammatica del sistema $\lambda 2$ o sistema F estende la grammatica del \ref{eq:simple_typed_lambda_grammar} tramite tre costrutti:
\emph{astrazione di termine su tipo}, \emph{applicazione di termine su tipo} \emph{quantificatore universale di tipo}. In tale sistema si permette a un termine di poter dipendere
da un tipo, tramite l'astrazione di tipo indicata con $\Lambda \alpha. M$. Il corrispondete costrutti dei tipi per l'astrazione di tipo è proprio il 
quantificatore universale, inoltre per poter $\beta$-ridurre tali costrutti è necessario l'applicazione di tipo che permettete di ridurre la corrispondente astrazione.
\begin{equation}
    \begin{array}{lrll}
    \text{Types:}   & \sigma, \tau & ::= \alpha                     & \text{(Variabili di tipo)}    \\
                    &              & \mid \sigma \to \tau           & \text{(Tipo funzione)}        \\
                    &              & \mid \forall \alpha. \sigma    & \text{(Tipo universale)}      \\
    \text{Terms:}   & M, N         & ::= x                          & \text{(Variabili)}            \\
                    &              & \mid \lambda x : \sigma.M      & \text{(Astrazione)}           \\
                    &              & \mid M N                       & \text{(Applicazione)}         \\
                    &              & \mid \Lambda \alpha. M         & \text{(Astrazione di tipo)}   \\
                    &              & \mid M \sigma                  & \text{(Applicazione di tipo)}
    \end{array}
    \tag{System F}
    \label{eq:grammar-system-f}
\end{equation}

\paragraph{Relazioni di calcolo e sostituzione.}
Le relazioni $\alpha$-equivalenza e $\eta$-conversione del sistema $\lambda 2$ rimangono invariate. Al 
contrario $\beta$-riduzione viene estesa con la medesima logica sui costrutti d'astrazione e d'applicazione di tipo. 

\begin{equation}
    (\Lambda \alpha. M) \tau \to_\beta M[\tau / \alpha]
    \label{eq:system_f_beta_reduction}
\end{equation}

\paragraph{Teorema di normalizzaione forte.}
Il sistema $\lambda 2$ gode della proprietà di normalizzazione forte, il che garantiche che l'esistenza di una forma normale
per ogni termine del sistema. Tale risultato fondamentale è stato dimostrato da Girard \cite{Girard1989}

\paragraph{Type checking.}
Per quanto riguarda il type checking nel sistema $\lambda 2$, e importante osservare che rimane decidibile. Il processo 
di verifica per il sistema $\lambda 2$ estende l'algoritmo \eqref{eq:simple_typed_lambda_type_cheking_algorithm} per 
gestire l'astrazione e l'applicazione di tipo, garantendo che le variabili di tipo siano gestite correttamente nel contesto.
\begin{equation}
    \text{type}(\Gamma, e) = 
    \begin{cases} 
        \sigma & \text{if } e = x \land (x : \sigma) \in \Gamma \\
        \sigma \to \text{type}(\Gamma, x:\sigma, e') & \text{if } e = \lambda x : \sigma . e' \\
        \forall \alpha. \text{type}(\Gamma, e') & \text{if } e = \Lambda \alpha. e' \land \alpha \notin FV(\Gamma) \\
        \sigma[\tau/\alpha] & \text{if } 
            \begin{cases} 
                e = e' \tau \\ 
                \text{type}(\Gamma, e') = \forall \alpha. \sigma 
            \end{cases} \\
        \tau & \text{if } 
            \begin{cases} 
                e = e_1 e_2 \\ 
                \text{type}(\Gamma, e_1) = \sigma \to \tau \\ 
                \text{type}(\Gamma, e_2) = \sigma 
            \end{cases} \\
        \perp & \text{otherwise}
    \end{cases}
\end{equation}

\paragraph{Type inference.} 
la type inference almeno nel caso generele è indecidibile, come dimostrato da Wells \cite{wells1994}. Per questo si utilizzano 
strategie basate su metavaribili e unificazione che pero nel caso peggiore richiederanno le annotazioni esplicite.

\paragraph{Esempi di termini e tipi.}
Nel sistema $\lambda 2$, il polimorfismo permette di definire termini che operano uniformemente su diversi tipi. 
Un esempio fondamentale è la funzione identità universale $ID \equiv \Lambda \alpha. \lambda x : \alpha. x$, 
il cui tipo è $\forall \alpha. \alpha \to \alpha$. Per applicare tale funzione a un valore concreto, è necessaria 
una fase preliminare di applicazione di tipo.
\begin{equation}
    (ID \: Nat \: 5 ) \to_\beta (\lambda x : Nat. x) \: 5 \to_\beta 5
\end{equation}
Questa espressività consente di codificare strutture dati logiche come i Booleani di Church, definiti dal tipo 
$\mathbb{B} \equiv \forall \alpha. \alpha \to \alpha \to \alpha$. In questa cornice, i valori di verità sono 
termini che selezionano il primo o il secondo argomento del tipo generico scelto.
\begin{align*}
    \mathbf{true}  &\equiv \Lambda \alpha. \lambda x:\alpha. \lambda y:\alpha. x \\
    \mathbf{false} &\equiv \Lambda \alpha. \lambda x:\alpha. \lambda y:\alpha. y
\end{align*}
L'applicazione $(\mathbf{true} \: Nat \: 1 \: 0)$ riduce a $1$ attraverso passi consecutivi di $\beta$-riduzione a 
livello di tipo e di termine, dimostrando come il polimorfismo gestisca l'astrazione sui dati in modo computazionale.
\subsection{Sistema \texorpdfstring{$\lambda\underline{\omega}$}{Lambda-omega-under}}

\paragraph{Grammatica.}
La grammatica del sistema $\lambda\underline{\omega}$ estende la grammatica del \ref{eq:simple_typed_lambda_grammar} tramite: una categoria sintattica chiamata kinds,  
il costrutto di \emph{astrazione di tipo su tipo} e \emph{applicazione di tipo su tipo}. In tale sistema si permette a un tipo di poter dipendere da un tipo, rendendo
i tipi delle funzioni, chiamate \emph{operatori di tipo}. Per garantire che gli operatori di tipo siano ben formate si utilizza la categoria sintattica kinds, ovvero 
il tipo dei tipi. Indroduzione dei costrutti di astrazione e applicazione di tipo su tipo traforma la categoria sintattica dei type in type constructors un calcolo 
separato sui tipi.
\begin{equation}
    \begin{array}{lrll}
        \text{Kinds:}               & K           & ::= \ast                & \text{(Kind terminale)}       \\
                                    &             & \mid K \to K            & \text{(Kind funzione)}        \\
        \text{Type Constructors:}   & A, B        & ::= \alpha              & \text{(Variabili di tipo)}    \\
                                    &             & \mid A \to B            & \text{(Tipo funzione)}        \\
                                    &             & \mid \lambda \alpha:K.A & \text{(Astrazione di tipo)}   \\
                                    &             & \mid A B                & \text{(Applicazione di tipo)} \\
        \text{Terms:}               & M, N        & ::= x                   & \text{(Variabili)}            \\
                                    &             & \mid \lambda x:A. M     & \text{(Astrazione)}           \\
                                    &             & \mid M N                & \text{(Applicazione)}
        \end{array}
    \tag{$\lambda\underline{\omega}$}
    \label{eq:grammar_system_omega_under}
\end{equation}

\paragraph{Relazioni di calcolo e sostituzione.}
Nel sistema $\lambda\underline{\omega}$, le relazioni di equivalenza si estendono per coprire la nuova sintassi degli operatori di tipo. 
La $\beta$-riduzione ($\to_\beta$) opera ora su due livelli, termini e tipi. Inoltre è importante notare che della dinamina la $\beta$-riduzione 
rimane invariata.

\begin{align}
    (\lambda x : \sigma. M) N & \to_\beta M[N/x] 
    \tag{Term-level $\beta$} 
    \label{eq:term_beta_reduction} \\
    (\lambda \alpha : K. \sigma) \tau & \to_\beta \sigma[\tau/\alpha] 
    \tag{Type-level $\beta$}
    \label{eq:type_system_omega_under_beta_reduction} 
\end{align}

\paragraph{Teorema di normalizzaione forte.}
Il sistema $\lambda\underline{\omega}$ essendo, come detto in precedenza, un doppio \ref{eq:simple_typed_lambda_grammar} gode della proprita di normalizzaione forte.
Il che garantiche che un termine abbia sempre una forma normale. Questo risultato è ottenibile mostrando che tra le categorie sintattiche addizionali di type constractor
e kinds ed il \ref{eq:simple_typed_lambda_grammar} esisteza un isomorfismo, ovvero possiedono comportameti indistinguibili sulla relazioni e riduzioni. Alla luce del 
risultato precedente si utilizza la stessa strategia con la categoria sintattica dei termini, considerando che i costruttori di tipo ammetono sempre forma normale.

\paragraph{Type checking.}
Il type checking del sistema, a differenza del calcolo semplicemente tipato (\ref{eq:simple_typed_lambda_grammar}) dove consiste in un mero confronto sintattico, 
acquisisce vere e proprie capacità di calcolo, percio di dice \emph{type checking modulo $\beta$-conversione}. Due termini appartengono alla stesso tipo sse la 
forma normale dei loro tipi è $\alpha$-equivalente. Cosiderando il terorema di normalizzazione forte a livello di costruttori di tipi ne segue la decidibilita.
Riportato di seguito il processo di verifica per il sistema $\lambda\underline{\omega}$ che estende l'algoritmo \eqref{eq:simple_typed_lambda_type_cheking_algorithm} 
per gestire i costruttori di tipo e i conseguenti kinds.
\begin{equation}
    \text{type}(\Gamma, e) =
    \begin{cases}
        A & \text{if } e = x \land (x : A) \in \Gamma \\
        A \to \text{type}(\Gamma, x:A, e') & \text{if } e = \lambda x : A . e' \land \text{kind}(\Gamma, A) = \ast \\
        B & \text{if }
            \begin{cases}
                e = e_1 e_2 \\
                \text{type}(\Gamma, e_1) \equiv_\beta A \to B \\
                \text{type}(\Gamma, e_2) \equiv_\beta A
            \end{cases} \\
        \perp & \text{otherwise}
    \end{cases}
\end{equation}

\paragraph{Type inference.}
% TODO

\paragraph{Esempi di termini e tipi.}
Nel sistema $\lambda\underline{\omega}$, l'introduzione dei costruttori di tipo permette di definire operatori che agiscono sulla struttura dei tipi stessi. 
Un esempio classico è l'operatore prodotto $\text{Pair}$, che può essere definito come un'astrazione di tipo su tipo con kind $\ast \to \ast \to \ast$.
\begin{equation}
    \text{Pair} \equiv \lambda \alpha : \ast. \lambda \beta : \ast. \alpha \to \beta \to \ast
\end{equation}
Si osservi che $\text{Pair}$ non è un tipo "abitabile" da termini, ma un operatore che richiede due tipi base per convergere a un tipo concreto ($\ast$). 
Applicando l'operatore ai tipi $Nat$ e $Bool$, otteniamo un nuovo tipo costruttore:
\begin{center}
    \setlength{\tabcolsep}{2pt}
    \begin{tabular}{r @{\hspace{5pt}} c @{\hspace{5pt}} l}
        $(\text{Pair} \: Nat \: Bool)$ & $\equiv$ & $(\lambda \alpha : \ast. \lambda \beta : \ast. \alpha \to \beta \to \ast) \: Nat \: Bool$ \\
        & $\to_\beta$ & $(\lambda \beta : \ast. Nat \to \beta \to \ast) \: Bool$ \\
        & $\to_\beta$ & $Nat \to Bool \to \ast$
    \end{tabular}
\end{center}
Si puo notare che il risultato del calcolo è un costruttore che riceve in input due tipi prefissati e restituise un'altro. È importante notare che la definizione 
del'operatore di tipo permettere di manipolare soltanto i tipi e non permetterano di scrivere funzioni che operano su termini.

\subsection{Sistema \texorpdfstring{$\lambda$}{Lambda}\texorpdfstring{$\Pi$}{Pi}}

\paragraph{Grammatica.}
La grammatica di $\lambda \Pi$ generalizza il costruttore di tipo $\to$ tramite il costruttore di \emph{prodotto dipendete} ($\Pi x : A.B$). Il prodotto dipendete permette di far dipendere un tipo
da un termine. Questa dipendeza è espressa dei costrutti di astarzione e applicazione di tipo su termine. Nello specifico la sintassi ($\Pi x : A.B$) indica una famiglia di tipi B indicizzata da
termini di tipo A. Si puo notare che ($A \to B$) non sia altro che un caso particolare in cui in $x \notin FV(B)$. Inoltre è importatene notare che si matengono i kinds con il medesimo costrutto
ma di livello superiore, tratta sempre la categoria sintattica dei termini.
\begin{equation}
    \begin{array}{lrll}
        \text{Kinds:}           & K           & ::= \ast                & \text{(Tipo dei tipi)}        \\
                                &             & \mid \Pi x : A. K       & \text{(Kind dipendente)}      \\
        \text{Types Families:}  & A, B        & ::= c                   & \text{(Costanti di tipo)}     \\
                                &             & \mid \Pi x : A. B       & \text{(Prodotto dipendente)}  \\
                                &             & \mid A M                & \text{(Applicazione di tipo)} \\
        \text{Terms:}           & M, N        & ::= x                   & \text{(Variabili)}            \\
                                &             & \mid \lambda x : A. M   & \text{(Astrazione)}           \\
                                &             & \mid M N                & \text{(Applicazione)}
    \end{array}
    \tag{$\lambda \Pi$}
    \label{eq:grammar_system_p}
\end{equation}

\paragraph{Relazioni di calcolo e sostituzione.}
Nel sistema $\lambda \Pi$ la distizione netta tra termini e tipi viene meno, poiche i tipi possono contenere termini, la struttura dei tipi stessa è soggeta al calcolo. 
Alla luce di ciò, si dirà che due tipi siano uguali se e solo se la riduzione del loro contenuto computazionale conduce a una forma normale comune. 
Questa relazione di equivalenza, denominata $\beta$-conversione ($A \equiv_\beta B$), implica che il type checking non sia più un semplice confronto sintattico, 
ma richieda l'esecuzione dei termini contenuti nei tipi stessi.
\begin{equation}
    (\Pi x : A. B) \, N \to_\beta B[N/x]
    \label{eq:lp_type_application}
\end{equation}
\begin{equation}
    (\Pi x : A. K) N \to_\beta K[N/x]
    \label{eq:lp_kind_application}
\end{equation}

\paragraph{Teorema di normalizzaione forte.}
Il sistema $\lambda \Pi$ gode della proprietà di normalizzaione forte, questo garantisce sempre l'esistenza di una forma normale.
Per una trattazione esaustiva del formalismo e della prova di normalizzazione per l'intero Cubo \cite{barendregt1992}.

\paragraph{Type checking.}
Nel sistema $\lambda \Pi$, il processo di type checking è decidibile ma richiede l'integrazione di fasi di calcolo esplicito. 
A differenza del calcolo semplicemente tipato, l'uguaglianza tra tipi non è un mero confronto sintattico, ma viene verificata 
modulo $\beta$-conversione e \eqref{eq:lp_type_application}. L'algoritmo di verifica deve quindi essere in grado di ridurre i 
termini contenuti nei tipi alle rispettive forme normali per accertarne l'equivalenza. 
\begin{equation}
    \text{gen}(\Gamma, e) =
    \begin{cases}
        (A, \emptyset) & \text{if } e = x \land (x : A) \in \Gamma \\
        (\Pi x : A. B, \mathcal{C}) & \text{if } e = \lambda x : A. e' \land (B, \mathcal{C}) = \text{gen}(\Gamma, x:A, e') \\
        (B[e_2/x], \mathcal{C}_{all}) & \text{if } e = e_1 e_2 \text{ dove:} \\ 
        & \begin{cases}
            (A', \mathcal{C}_1) = \text{gen}(\Gamma, e_1)   \\
            (A'', \mathcal{C}_2) = \text{gen}(\Gamma, e_2)  \\
            \mathcal{C}_{all} = \mathcal{C}_1 \cup \mathcal{C}_2 \cup \{A' \equiv_\beta \Pi x : A. B, A'' \equiv_\beta A\}
        \end{cases}
    \end{cases}
\end{equation}

\paragraph{Type inference.}
Nel sistema $\lambda\Pi$, il processo di \emph{type inference} rimane decidibile. Data un'applicazione $(M \, N)$, l'algoritmo determina 
il tipo del risultato non solo verificando la coerenza tra il tipo dell'argomento e il dominio della funzione, ma effettuando una sostituzione 
esplicita. Il tipo finale viene ottenuto sostituendo il termine $N$ all'interno del prodotto dipendente $\Pi x:A. B$, come formalizzato in 
\eqref{eq:lp_type_application}. Questo passaggio è cruciale, poiché il tipo dell'output è ora funzione del valore dell'input.

\paragraph{Esempi di termini e tipi.}
L'esempio classico di termine per i tipi dipendenti è il vettore di dimensione $n$. Si cosideri una famiglia di tipi $Vect : Nat \to \ast$. 
Possiamo definire su tale famiglia una funzione $\text{head}$ che estrae il primo elemento di un vettore di lunghezza $n+1$, garantendo staticamente 
l'assenza di errori out of bounds.
\begin{equation}
    \text{head} : \Pi n : Nat. VectA \, (n+1) \to A
\end{equation}
È importante notare che A deve essere un tipi prefissato, poiche nel sistema $\lambda\Pi$ non è presente polimorfismo. Se tentassimo di applicare 
$\text{head}$ a un vettore di tipo $Vect \, 0$, l'algoritmo di unificazione fallirebbe poiché $0$ non è unificabile con $n+1$ per alcun $n \in Nat$. 
Il type checking diventa così una forma di verifica formale della correttezza del software.

\subsection{Combinazioni}
Una volta introdotti tutti i tre assi è possibile osservare, come detto precedentemente, che esisteno otto combinazione di sistema di tipo distinti.
Ciascuno di essi aggiunge una specifica capacità di astrazione al calcolo semplicemente tipato ($\lambda^\to$), portando alla gerarchia sintetizzata 
nella Tabella \ref{tab:barendregt_combinations}. La medesima gerarchia sara fondamentale per definire il ponte con tutti i sistemi logici. I nomi dei 
sistemi di tipo per la maggior parte delle combinazioni sono l'unione dei nomi dei sistemi uniti.

\begin{table}[H]
    \centering
    \small
    \begin{tabular}{ccccc}
        \hline
        \textbf{Sistema}                 & \textbf{Term-Term} & \textbf{Term-Type} & \textbf{Type-Type}             & \textbf{Type-Term} \\ 
                                         & ($\lambda^\to$)    & ($\lambda 2$)      & ($\lambda \underline{\omega}$) & ($\lambda \Pi$)    \\ \hline
        $\lambda^\to$                    & \checkmark         & --                 & --                             & --                 \\
        $\lambda 2$                      & \checkmark         & \checkmark         & --                             & --                 \\
        $\lambda \underline{\omega}$     & \checkmark         & --                 & \checkmark                     & --                 \\
        $\lambda \omega$                 & \checkmark         & \checkmark         & \checkmark                     & --                 \\
        $\lambda \Pi$                    & \checkmark         & --                 & --                             & \checkmark         \\
        $\lambda \Pi 2$                  & \checkmark         & \checkmark         & --                             & \checkmark         \\
        $\lambda \Pi \underline{\omega}$ & \checkmark         & --                 & \checkmark                     & \checkmark         \\
        CoC                              & \checkmark         & \checkmark         & \checkmark                     & \checkmark         \\ \hline
    \end{tabular}
    \caption{Le otto combinazioni del cubo di Barendregt.}
    \label{tab:barendregt_combinations}
\end{table}

\section{Calcolo del costruzioni (CoC)}
\label{sec:coc}
Il calcolo delle costruzioni è l'apice dell'espressivita dei sistemi di tipo. Per ottenere questa potenza espressiva si unificando le grammatiche di tutte le
direzioni del cubo precedentemente discusso, eliminando tutti i costrutti ridondanti. Nei sistemi predendeti si dividevando nettamete le categorie 
sitattiche: termini, tipi e kind. Questa divisione portava ad avere costrutti uguali ma separati per categorie sintattiche, l'esempio migliore l'astrazione e l'applicazione.
Inoltre è importante notare che i costrutti vengono unificati, non persi, percio lo studio dei sistemi precendeti rimande fondamentale per la compresione e l'attibuzione ai
costrutti del potere espressivo. Quindi spesso ci utilizzaremo nomenglature dei precedeti sistemi.

\subsection{Grammatica}
La grammatica del calcolo delle costruzioni, come anticipato, molto più ridotta in numero di costrutti rispetto alla fusione cieca dei sistemi. Nel calcolo delle costruzioni
è presente solo una categorie sintatticha, divisa per espressivita in: sort e termini. Questa unificazione è il presupposto logico che permette al prodotto dipendente ($\Pi$) 
di agire come un costruttore universale. Grazie alla possibilità di quantificare indistintamente su basi atomiche ($\ast$) o su universi superiori ($\square$), il tipo $\Pi$ 
è in grado di sussumere e generalizzare tutti i costrutti caratterizzanti dei sistemi precedenti, quali l'implicazione funzionale ($A \to B$), il quantificatore universale del 
secondo ordine ($\forall \alpha : \ast . T$), gli operatori di tipo e i tipi dipendenti propriamente detti.

\begin{equation}
    \begin{array}{lrll}
        \text{Sorts:}           & s           & ::= \ast \mid \square   &                               \\
        \text{Pseudo-terms:}    & A, B, M, N  & ::= x                   & \text{(Variabili)}            \\
                                &             & \mid s                  & \text{(Costanti)}             \\
                                &             & \mid \lambda x : A.M    & \text{(Astrazione)}           \\
                                &             & \mid \Pi x : A.B        & \text{(Prodotto dipendente)}  \\
                                &             & \mid M N                & \text{(Applicazione)}
    \end{array}
    \tag{CoC}
    \label{eq:coc_grammar}
\end{equation}

\subsection{Relazioni di calcolo e sostituzione}
Il Calcolo delle Costruzioni non introduce nuove regole di riduzione rispetto al $\lambda$-calcolo puro; di fatto, permangono valide tutte le definizioni di 
$\alpha$-conversione, $\beta$-riduzione ed $\eta$-conversione analizzate nei capitoli precedenti. Tuttavia, la vera innovazione risiede nell'estensione del 
loro dominio di applicazione.Mentre nei sistemi più semplici (come il sistema $\lambda{\to}$) le riduzioni operano esclusivamente sul piano dei termini (i programmi),
nel CoC esse agiscono in modo uniforme sull'intera categoria dei pseudo-termini. Ciò significa che la computazione non è più confinata alla fase di esecuzione, 
ma diventa parte integrante della definizione stessa dei tipi.

\subsection{Unificazione dei costrutti nel CoC}
L'unione delle categorie sintattiche permette al prodotto dipendente ($\Pi x : A. B$) di agire come \emph{costruttore universale}. L'unificazione dal CoC permette 
al tipo $\Pi$ di assumere significati logici e computazionali differenti a seconda della natura (ovvero del \textit{sort}) dei termini coinvolti:
\begin{description}
    \item \textbf{Implicazione funzionale ($A \to B$):}
        Se $A : \ast$, $B : \ast$ e la variabile $x$ non compare libera in $B$ ($x \notin FV(B)$), 
        il prodotto dipendente degenera nella classica freccia del lambda calcolo semplicemente tipato ($\lambda{\to}$).
    \item \textbf{Polimorfismo ($\forall \alpha : \ast . T$):} 
        Se $A : \square$ (un genere) e $B : \ast$ (un tipo), il $\Pi$ permette di quantificare sopra un intero universo di tipi. 
        Questo unifica il CoC con il sistema $\lambda2$, permettendo la definizione di funzioni generiche che accettano tipi come argomenti.
    \item \textbf{Operatori di tipo:}
        Se sia $A$ che $B$ appartengono a $\square$, il prodotto definisce funzioni tra generi, 
        permettendo la creazione di costruttori di tipi complessi (come \texttt{List} o \texttt{Tree}) 
        tipici del sistema $\lambda\underline{\omega}$.
    \item \textbf{Tipi dipendenti ($Vect(n)$):}
        Se $A : \ast$ (un tipo di dato) e $B : \square$ (un genere), il tipo risultante dipende dal valore di un termine. 
        È questo il cuore della logica dei predicati, dove una proprietà $B$ è parametrizzata da un elemento $x$ di $A$.
\end{description}

\begin{comment}
    \subsection{Teorema di normalizzaione forte}
    % TODO
    \subsection{Type checking e Type inference}
    % TODO
    \paragraph{Esempi di termini e tipi.}
    % TODO
\end{comment}